% A Priori Knowledge of the World: Kant's Question and Three Computed Answers
% Claes Johnson.  DRAFT.
% Companion to math_education_after_ai.tex (which argues that solving is not empty),
% ai_vs_numerical_analysis.tex and constructive_turbulent_ns.tex.  This one takes the
% same thesis -- latent information made explicit by labour -- and puts it where it
% belongs, as an answer to Kant's question about synthetic a priori knowledge, with
% three worked examples of rising difficulty: the circle, the Euler equations, RealQM.
\documentclass[11pt]{article}
\usepackage[margin=1.15in]{geometry}
\usepackage[T1]{fontenc}
\usepackage{amsmath,amssymb}
\usepackage{url}
\usepackage[hidelinks]{hyperref}
\setlength{\parskip}{3pt}

\title{\vspace{-2em}A Priori Knowledge of the World:\\Kant's Question and Three Computed Answers}
\author{Claes Johnson\\[0.3em]
\small KTH Royal Institute of Technology, Stockholm\\
\small \texttt{claesjohnson@gmail.com}}
\date{\today}

\begin{document}
\maketitle

\begin{abstract}
\noindent
Kant asked how synthetic a priori knowledge is possible: knowledge which says something about the
world and yet owes nothing to observation. His own answer, that the knower supplies the form of
experience, is not believed any more, and his leading example --- that physical space is Euclidean,
knowably so in advance --- is taken to have been refuted. The question survived the answer, and it
is now answerable, because the stock of examples has grown in a direction he could not have
anticipated: the a priori knowledge arrives out of a computation. I give three, of rising
difficulty, in which a quantity belonging to the world is obtained with no measurement of it: the
circumference of the circle $x^2+y^2=1$; the drag and lift of a body in incompressible slightly
viscous flow, from the shape of the body alone; and the configuration of a molecule, from its atoms
alone.

What makes all three a priori is not that no physical process occurred but that \emph{the object was
never interrogated} --- no circle was drawn, no wing placed in a tunnel, no water in the room. A
spectrometer interrogates the molecule and its response makes the number; a simulation consults no
water, and would give the same digits in a universe containing none. The answer to Kant's question
is then available in a form he could not have reached: such knowledge is \emph{latent information
made explicit by labour}. The content was in the input, so nothing was created; but latent is not
available, and the only route from implicit to explicit is a chain of steps. That both deflates the
mystery and relocates it. The question is no longer whether the knowledge is possible but at what
rate the implicit becomes explicit and whether the rate can be certified --- exact for the circle, a
computed number with a bound for the other two, and what separates them is dissipation.

The a priori then goes all the way down, further than is usually conceded. No experience enters the
definition of the circle; nor Newton's second law and mass conservation, which are what the Euler
equations consist of, and whose one measured parameter drops out of the answer; nor the inverse
square law, whose form is forced by a conserved flux in three dimensions rather than fitted to a
torsion balance. RealQM's whole input is a geometric law and a list of integers. Experience enters
at one point only, in the claim that a construction applies to some given object --- and that claim
is the antecedent of a conditional, which is to say the specification, supplied rather than
discovered. Nobody calls the circumference of a circle empirical because somebody had to choose the
radius.

That no fitted quantity appears anywhere in the three is moreover a requirement and not an accident
of tidy construction, and the requirement is Leibniz's. An arbitrary parameter is exactly a quantity
for which a theory can give no reason, so by the principle of sufficient reason it cannot stand:
either the theory determines it, and it was never arbitrary, or the theory is unfinished at that
point. A fitted parameter is not a measurement incorporated into a theory but a hole in the theory
with a number in it --- which is the charge against turbulence constants, against functionals
calibrated on a training set whose answers were known in advance, and against contracted basis sets.
The same principle disposes of the two concessions ordinarily made to Kant's critics. Three
dimensions is not an empirical input but forced, since below three nothing can be complicated and
above three nothing can be stable, which leaves one thinkable option; and a curvature adopted
without a reason for that curvature rather than another is an arbitrary parameter in geometrical
costume.
\end{abstract}

\section{Kant's question}\label{sec:question}

The \emph{Critique of Pure Reason} is organised around a question which is easy to state and hard to
dismiss. Some judgements are analytic: the predicate is already contained in the subject, as in
\emph{bachelors are unmarried}. These are certain and they are empty, telling us about our words
rather than about the world. Other judgements are synthetic a posteriori: they add something, and
they are bought with experience, as in \emph{this rod is one metre long}. Between the two Kant
placed a third class, synthetic a priori: judgements which add something about the world and are
nonetheless known without consulting it~\cite{kant,prolegomena}. He held that mathematics is of this kind, and the
principles of pure natural science with it, and he asked how that is possible.

His answer was the one thing in the \emph{Critique} everybody remembers. It is possible, he said,
because the form of experience is contributed by the knower rather than found in the object. Space
and time are forms of sensible intuition; substance and causality are forms of the understanding. We
can know in advance what must hold of any possible object of experience because we supply the mould
into which experience is poured. Objects conform to our mode of cognition, not the reverse.

That answer is not believed now, and Kant's own leading instance of it failed: he held the geometry
of physical space to be Euclidean, knowably so a priori, and it is not. The question nonetheless did
not go away, and in the two and a half centuries since it was asked the stock of examples has grown
in a direction Kant could not have anticipated. The a priori knowledge now arrives out of a
computation. I give three instances in order of rising difficulty and then ask what they have in
common.

\section{First example: the circle}\label{sec:circle}

Define a curve by $x^2 + y^2 = 1$. Its circumference is $2\pi = 6.283185\ldots$, and one may have as
many further digits as one is willing to pay for. Nothing was measured. No circle was drawn, no
string laid along a rim, no instrument consulted. A quantity belonging to a figure in space was
obtained from a definition.

Two features of this are worth recording, because both recur in the harder examples.

\paragraph{The definition does not contain the digits.} It is a short formula. The number it
determines is transcendental --- Lindemann settled that in 1882~\cite{lindemann} --- so there is no
finite expression for it in rationals and no amount of staring at $x^2+y^2=1$ will produce
$3.14159$. The digits are obtained by \emph{construction}: Archimedes, in
\emph{Measurement of a Circle}~\cite{archimedes}, inscribed and circumscribed regular polygons of
ninety-six sides and squeezed the value between $223/71$ and $22/7$ --- a procedure which must be
run, and which yields as many further digits as one is prepared to double the number of sides for.
Every later method is likewise a procedure which must be run. So the knowledge is a priori and it is
also work. \emph{A priori does not mean free.} It names the
source of the knowledge, not its cost.

\paragraph{The knowledge is conditional.} The claim is not that any round thing in the world has
circumference $2\pi$ times its radius. It is that \emph{if} the shape is the one defined, \emph{then}
the circumference is that. The condition is not a weakening. It is what makes the statement exact,
and it is the form in which the a priori survived the collapse of Kant's version of it
(\S\ref{sec:kant-wrong}).

\section{Second example: drag and lift from the shape of the body}\label{sec:euler}

Aerodynamics was empirical in its foundations. The drag coefficient of a sphere, a cylinder, an
aerofoil is a tabulated quantity, and the tables came from wind tunnels. An engineer sizing a
structure looks up a number which somebody else measured on a model of something like it.

The claim that this detour through measurement is unnecessary is not a modern one. Euler
formulated his equations in 1755~\cite{euler1755} and claimed for them exactly that: that they
contain all of slightly viscous, nearly incompressible flow --- water, and air at subsonic speeds
--- and that nothing further need be supplied. The claim was made and then not acted on, because
there was no way to solve the equations. What the computational programme behind this article adds
--- RealEuler, the fluid counterpart of RealQM --- is the solving, and with it the vindication of
Euler's claim~\cite{cfd,dalembert,flight}.

Solve the Euler equations for incompressible flow with slight viscosity, with a slip or small
friction condition at the body, and the drag and the lift come out of the shape of the body alone.
No turbulence model fitted to data, no measured coefficient, no boundary layer resolved, no
circulation inserted by hand to make a wing lift. The input is a shape. The output is a force.

What had to be cleared out of the way first was d'Alembert's paradox, which stood for two centuries:
the inviscid incompressible equations admit potential flow, and potential flow has zero drag, so the
equations appeared to contradict the most elementary observation about fluids. The resolution is that
potential flow is unstable~\cite{dalembert}. It is a solution and it is not the one which occurs;
what occurs is a turbulent wake whose formation is itself computable, and which carries the drag.
The lift of a wing comes out of the same computation, with no circulation
inserted~\cite{flight}.

The honest form of the claim is already in the record of the programme, and it is stranger than the
claim itself: mean drag and lift converge to a few per cent on meshes on which the velocity field
does not converge and never will~\cite{drag,parameterfree,clay}. The \emph{output} is knowable while the \emph{solution} is not. I
return to this in \S\ref{sec:form}, because it is the clearest sign that what is being computed is
not the fluid but a functional of it.

The two features of \S\ref{sec:circle} reappear. The knowledge is a priori and it is expensive: a
mesh, a solver, a stopping criterion, and a number which arrives with a bound rather than exactly.
And it is conditional: given this shape, that force.

\section{Third example: the configuration of a molecule from its atoms}\label{sec:realqm}

Give the nuclear charges and how many electrons there are. Nothing else. No measured bond lengths,
no force constants taken from spectroscopy, no basis set contracted against molecular data, no
exchange--correlation functional calibrated on a training set of molecules whose answers were known
in advance. Then minimise the energy of a collection of non-overlapping charge densities, one per
electron, interacting by Coulomb's law with each pair counted once~\cite{realqm,realqmbook}. The configuration of the
molecule is what comes out.

From the record~\cite{universality}: for water, the equilibrium O--H bond comes out at
$1.814\,a_0$ --- which is the geometry the calculation is \emph{started} from rather than a relaxed output, so it should be read as an input --- and the H--O--H angle, which \emph{is} an output, at $103.3^\circ$ against the measured $104.5^\circ$. The lone pair is
not put in by hand --- the sector of the oxygen valence pointing away from both hydrogens is simply
where the energy puts it. Four free densities about a single kernel, seeded lopsidedly, relax into
the tetrahedral basin rather than the square one~\cite{chirality}, which is methane's
$109.47^\circ$ arrangement arrived at rather than assumed.

What does \emph{not} come out a priori is the absolute energy. The programme's own verdict on itself
is that the geometry is right and the energies are wrong~\cite{atoms}: errors cancel in differences taken within
one composition and do not cancel across compositions. This is the most interesting of the three
cases rather than the weakest, and \S\ref{sec:form} is about why.

\section{What makes the knowledge a priori}\label{sec:criterion}

An objection arrives immediately. A computer is a physical object; a computation is a physical
process, taking time, dissipating heat, occasionally giving the wrong answer because a transistor
was struck by a cosmic ray. If the number came out of a physical process, in what sense was it
obtained without observation?

The criterion is not whether some physical process occurred. It is whether \emph{the object was
interrogated}. A spectrometer interrogates the molecule: light goes in, the molecule responds, and
the response makes the number, so the molecule is a cause of the result. A simulation of water
consults no water. The silicon is physical and it stands in no causal relation whatever to any
sample of H$_2$O; one could run it in a universe containing no water at all and get the same
digits. In the three examples no circle was drawn, no wing was placed in a tunnel, and there was no
water in the room.

This is a sharper criterion than ``independent of experience'', which is vague enough to have
absorbed two centuries of argument. It also has a clean consequence: the a priori character of a
result is a fact about its causal provenance, and can therefore be settled by asking what was in the
room.

\section{How it is possible}\label{sec:how}

The deflationary half of the answer should be stated first, because it is true. In each case the
answer was already in the input. The definition of the circle determines $2\pi$. The Euler equations
together with the shape determine the drag. Coulomb's law together with the non-overlap condition
determines the geometry. The computation creates no information about the world; it transfers
information from the law to the prediction.

The other half is equally true and it is the whole point. \textbf{Latent is not available.} Nobody
can read the H--O--H angle off Coulomb's law by inspection, and nobody ever read $\pi$ off
$x^2+y^2=1$. The content is implicit, and the only route from implicit to explicit is a chain of
steps. This is the argument made at length elsewhere about why solving is not
empty~\cite{education,ainumerical}: the equation posits exactly as axioms do, its solution is implied but unavailable, and a chain of steps ---
Newton's method as readily as a proof --- constructs what was not there before. Nobody says that a
computed solution of a differential equation is empty on the ground that the equation already
determined it.

So the answer to Kant's question is:

\begin{quote}
A priori knowledge of the world is \emph{latent information made explicit by labour}.
\end{quote}

The possibility was never the mysterious part. What was missing from the old debate was any account
of the labour, and that is exactly what was unavailable before there were machines to do it, because
the labour in the interesting cases is far beyond hand calculation.

Once it is put this way, Kant's question changes shape rather than getting answered on its own
terms. He asked whether such knowledge is possible. The computational era answers yes, routinely,
industrially, and replaces the question with a better one: \emph{at what rate does the implicit
become explicit, and can the rate be certified?} For $\pi$ the cost is known in advance and the
digits are exact. For drag the cost is a mesh and the certificate is a bound obtained by duality
against a computed adjoint solution. For a molecular geometry the certificate is a mesh ladder and a
step ladder, and a difference below the drift of the iteration carries no information. The modulus
is the content --- which is the same conclusion one reaches coming at this from the side of
constructive existence, and it is reassuring that two routes arrive at it~\cite{constructive}.

\section{How far down does the a priori go?}\label{sec:kernel}

It is natural to concede at this point that the three examples are a priori only above an empirical
foundation --- that observation bought the laws, and that computation merely spends them. The
concession is too quick, and it dissolves on inspecting what actually enters.

\paragraph{The circle.} Nothing empirical enters. $x^2+y^2=1$ is a stipulation, and the derivation
of $2\pi$ from it consults nothing further.

\paragraph{Euler.} The equations are momentum balance for a continuum, which is to say Newton's
second law together with conservation of mass and the kinematic constraint of incompressibility.
That is the whole of their content, and it is why Euler could claim in 1755 that they contain all of
slightly viscous nearly incompressible flow: there was no further law to add, and no coefficient to
go and measure.
Whether Newton's second law is an empirical discovery at all is a long-standing question, and the
answer of Mach and Poincar\'e is that it is not~\cite{mach,poincare}: force is \emph{defined} by
$ma$, and the law is a
convention which fixes the meaning of its terms rather than a generalisation someone observed to
hold. Nothing in the chain is of the form ``and we measured it to be so''. Moreover the one
dimensional parameter which does carry a measured value, the viscosity, drops out of the answer: the
drag coefficient of a bluff body is essentially independent of Reynolds number over the whole range
of interest~\cite{drag}, which is precisely why the shape alone suffices.

\paragraph{RealQM.} The inverse square law is what a conserved radial flux in three dimensions must
look like. Its form is forced by Gauss's law and the dimension of space, not fitted to Coulomb's
torsion balance --- the balance confirmed a law whose shape was not open. The coupling strength is a
choice of units, and in atomic units it is one. What remains to be supplied is the nuclear charges
and the number of electrons, and those are integers. So the entire input to a RealQM calculation is
a geometric law and a list of integers. There is no measured number in it.

Then where does experience enter? In one place only: in the claim that these constructions apply to
the actual world --- that physical space is such that the defined circle is realisable in it, that
air at these speeds behaves as an incompressible continuum, that matter consists of nuclei and
electrons in these numbers. And by \S\ref{sec:kant-wrong} that claim is no part of the knowledge.
It is the antecedent of a conditional, which is to say it is the \emph{specification}, and a
specification is supplied rather than discovered. ``Given these nuclei'' stands exactly where
``given this shape'' stands in the first example, and nobody calls the circumference of the circle
empirical on the ground that somebody had to choose the radius.

The right conclusion is therefore the stronger one. \textbf{The a priori goes all the way down.}
Experience enters not as a premise of any derivation but at the moment of application, when one
decides to point a construction at a particular object; and that decision is not knowledge about the
world but a choice about what to compute.

One might object that the inverse square law depends on the dimension of space being three, and that
three is a fact about this world rather than a theorem. It is not an input either, and
\S\ref{sec:kant-wrong} is the reason: below three nothing can be complicated and above three
nothing can be stable, so the dimension is forced by the requirement that there be matter at all.

What does remain is one thing, and it is worth naming because it falls exactly where the examples
are weakest. Absolute scale must come from somewhere. Atomic units fix it by convention, so a
geometry in Bohr radii is unit-free while an energy in Hartree is not comparable with a measured one
until the scale is pinned. That residue attaches to \emph{magnitude} rather than to \emph{form},
which is the distinction \S\ref{sec:form} reaches from the other direction. The two meeting is not
an accident: where the empirical residue sits is where the programme's quantitative weakness sits
too.

\paragraph{Why there could not have been a measured number.} The absence of fitted quantities in the
three examples is not a matter of taste, or of having built tidily. It is required, and the
requirement is Leibniz's.

The principle of sufficient reason holds that nothing is so without a reason why it is so rather
than otherwise. Leibniz used it destructively and to great effect. If space were absolute, there
could be no sufficient reason for the world to sit here rather than a yard to the left --- so, he
concluded against Clarke, space is not absolute~\cite{leibnizclarke}. The same argument disposed of absolute time: no
reason could be given for the world to have begun at this moment rather than six months earlier. It
is the \emph{lack} of a sufficient reason that does the work, and the thing which lacks one is
abolished.

Turn that on a parameter. An arbitrary parameter is precisely a quantity for which the theory
supplies no reason: nothing within it says why the value is this rather than that, which is why the
value has to be fetched from outside and fitted. By the principle such a quantity cannot stand.
Either the theory determines it, in which case it was never arbitrary; or the theory offers no
reason for it, in which case the theory is unfinished at that point. A fitted parameter is not a
measurement incorporated into a theory. It is a hole in the theory with a number in it.

So \textbf{there can be no arbitrary parameters}, and the three examples contain none --- not by
good fortune, but because anything else would be defective. This also says, in one word, what is
wrong with the alternatives they displace. A turbulence model carries constants fitted to measured
flows, which is what makes a parameter-free method worth the name~\cite{parameterfree}. A density functional carries parameters calibrated against a training set of molecules whose
answers were known in advance. A contracted basis set is fitted. In each case the fitting marks the
place where the theory stopped explaining and began recording, and each constant is a point at which
it has no sufficient reason to offer. That such a theory then reproduces the data it was fitted to
is not evidence in its favour; it is what having been fitted means.

The principle is a demand rather than a theorem, and the honest statement of its force is that it
convicts a theory of being unfinished and not of being false --- a model with free parameters can
be perfectly serviceable, and most of physics has been. But unfinished is the right verdict, and it
is the verdict which says what remains to be done.

\section{Three dimensions, Euclidean, and no assumption required}\label{sec:kant-wrong}

Kant's flagship case is usually said to have failed. He held Euclidean geometry to be synthetic a
priori and true of physical space; Gauss, Bolyai and Lobachevsky produced consistent alternatives,
Riemann gave the general framework in 1854~\cite{riemann}, general relativity arrived in 1915, and the geometry of
space became a contingent matter to be settled by observation. That is taken to end the project.

The refutation is weaker than advertised, and it is worth separating the two things it is supposed
to have established.

\paragraph{Three is not one option among many.} It is the only dimension in which there is anything
to compute, and the argument for that is an argument and not a measurement.

Below three, there is no complexity to be had. In two dimensions a network of connections must be
planar, since nothing can cross without meeting, so no circuit of any depth can be wired; there are
no knots; and a passage through an organism from one side to the other severs it into two. Whitrow
put this case in 1955~\cite{whitrow} and it has not been answered; the later literature surveys it
without improving on it~\cite{barrowtipler}.

Above three, there is nothing stable to compute \emph{with}. A conserved flux in $d$ dimensions
gives a potential going as $r^{-(d-2)}$ and a force as $r^{-(d-1)}$, and the stability of orbits
under such a force fails at $d = 4$: by Bertrand's theorem~\cite{bertrand} only the inverse square and the linear
law return closed orbits at all, and the inverse cube force of four-dimensional space gives
marginal instability and then, above it, spiralling. The quantum version is worse and more decisive:
the four-dimensional Coulomb potential goes as $r^{-2}$, which is exactly the critical case for
collapse to the centre, so the atom has no ground state to sit in. This is Ehrenfest's argument of
1917~\cite{ehrenfest}, and it says that in four dimensions there are no atoms, hence no chemistry, hence nothing for
a theory of matter to be about. Signalling fails in the same place: sharp wave propagation ---
Huygens' principle --- holds in odd dimensions from three upward and fails in even ones, so a pulse
in two or four dimensions drags a tail behind it and carries no clean information.

So the dimension three is not an empirical input which observation happened to supply. It is forced,
and forced by the requirement that there be bound matter and complexity at all. Below it nothing can
be complicated; above it nothing can be stable. That leaves one thinkable option, and the
conditional form in which this holds --- if there is to be matter and complexity, the dimension is
three --- is the same form as every other claim in this article.

\paragraph{Euclidean needs no reason; the alternative does.} Riemann's generalisation is
mathematically free, and mathematical freedom is not physical motivation. Consistency establishes
that non-Euclidean geometries can be \emph{written down}, not that there is any ground for
positing one, and the burden lies on whoever does. This is the argument of
\S\ref{sec:kernel} once more, and it is the one Leibniz aimed at absolute space: a curvature
adopted without a reason for that curvature rather than another is an arbitrary parameter wearing
a geometrical costume. The ground usually offered is gravitation, and
the programme of which this article is a part does not take gravitation to be geometry: charge and
energy generate it, and it emerges above the atomic scale rather than being built into a
metric~\cite{gravity}. On
that account curvature is a way of describing a dynamics and not a property of space, and space
stays flat. That is a commitment of the programme rather than a result of it, and it is recorded
here as such --- but it means the inference from gravitation to curved space is a choice, and one
this programme declines.

\paragraph{And in any case, the conditional form.} Whether or not the two paragraphs above are
granted, none of the three examples asserts a geometry of the world unconditionally, so none repeats
the mistake Kant is accused of. Each says: given this specification, that consequence. The circle
example does not claim that physical space is Euclidean; it claims that the curve so defined has
that circumference. The flow example does not claim that air is an incompressible continuum; it
claims that a body of this shape in such a fluid feels that force. RealQM does not claim that the
inverse square law is a necessary truth; it claims that charge obeying it, so constrained, arranges
itself so. In conditional form the a priori is not merely defensible but ordinary. It is what every
computation does.

\section{Form before magnitude}\label{sec:form}

The three examples do not behave alike, and the difference is instructive rather than embarrassing.
It is a ladder of rising irreducibility.

The circle has a closed form, so there is nothing to certify: $2\pi$ is the answer and the only
question left is how many digits one wants. The flow has no closed form, and what can be certified
is not the flow but a functional of it --- the drag arrives with a bound obtained by duality, and it
arrives on meshes which do not resolve the velocity field and never will. The molecule has no closed
form either, and here the certificate narrows once more: it is trustworthy for differences taken
within one composition, where the errors cancel, and not for absolute magnitudes, where they do not.

What rises along that ladder is irreducibility of the computation. And what the second and third share,
which the first does not, is \emph{dissipation}. The flow sends energy to scales no mesh resolves,
and the drag survives that only because it is a mean value. The molecular configuration is reached
by a relaxation which approaches equilibrium asymptotically and never arrives, so that every
quantity read off it is read off at a finite time and carries the drift of the iteration. Where
there is dissipation there is no closed form, and the a priori must then be bought with a modulus
rather than stated in a formula.

The pattern in what RealQM gets right and wrong is then sharp. It delivers \emph{form} a priori ---
which site, which branch, which geometry, which symmetry, tetrahedral rather than square, where the
lone pair points --- and not \emph{magnitude}. The structural results do not depend on the energies
being right, which is why they survive the admission that the energies are not.

There is an odd echo of Kant in this, who also held that what the a priori supplies is form. I do
not want to press a resemblance into an argument, and the sense of ``form'' is not his. But it is
worth noticing that when a programme sets out to compute the world from its law and nothing else,
the first thing to come out is shape.

\section{What it changes}\label{sec:changes}

If drag follows from shape and molecular geometry follows from charge, then the wind tunnel and the
diffractometer change their role. They stop being where the number comes from and become where the
theory is tested. That is not a demotion --- a theory with no check against the world is worth
nothing, and the three examples are credible only because their outputs have been compared with
measurement. But it reverses the direction of supply, and that reversal is the practical content of
the whole argument.

It also places the frontier for automation exactly where the companion articles place it. The labour
by which the implicit becomes explicit is machine work now, and increasingly good machine work;
systems have begun to produce research-level proofs and they write and run solvers. What is not machine
work is the formulation and the decision about what the object is ---
and not because the technology is young. There is nothing to derive a formulation from, so no chain
of steps produces one.

Kant asked how we can know anything about the world without looking at it. The answer, two and a
half centuries on, is that we look once, carefully, and then compute.

\begin{thebibliography}{99}

\bibitem{kant} I.~Kant, \emph{Kritik der reinen Vernunft}, Riga, 1781; 2nd edn.\ 1787.

\bibitem{prolegomena} I.~Kant, \emph{Prolegomena zu einer jeden k\"unftigen Metaphysik}, Riga, 1783.

\bibitem{leibnizclarke} G.~W.~Leibniz and S.~Clarke, \emph{The Leibniz--Clarke Correspondence},
1715--1716; ed.\ H.~G.~Alexander, Manchester University Press, 1956.

\bibitem{archimedes} Archimedes, \emph{Measurement of a Circle}, Proposition~3; in T.~L.~Heath
(ed.), \emph{The Works of Archimedes}, Cambridge University Press, 1897.

\bibitem{lindemann} F.~Lindemann, \emph{\"Uber die Zahl $\pi$}, Math.\ Ann.\ \textbf{20} (1882)
213--225.

\bibitem{mach} E.~Mach, \emph{Die Mechanik in ihrer Entwicklung historisch-kritisch dargestellt},
Brockhaus, Leipzig, 1883.

\bibitem{poincare} H.~Poincar\'e, \emph{La Science et l'Hypoth\`ese}, Flammarion, Paris, 1902.

\bibitem{riemann} B.~Riemann, \emph{\"Uber die Hypothesen, welche der Geometrie zu Grunde liegen},
Habilitationsvortrag, G\"ottingen, 1854.

\bibitem{whitrow} G.~J.~Whitrow, \emph{Why physical space has three dimensions},
Br.\ J.\ Philos.\ Sci.\ \textbf{6} (1955) 13--31.

\bibitem{ehrenfest} P.~Ehrenfest, \emph{In what way does it become manifest in the fundamental laws
of physics that space has three dimensions?}, Proc.\ Amsterdam Acad.\ \textbf{20} (1917) 200--209.

\bibitem{bertrand} J.~Bertrand, \emph{Th\'eor\`eme relatif au mouvement d'un point attir\'e vers un
centre fixe}, C.~R.\ Acad.\ Sci.\ Paris \textbf{77} (1873) 849--853.

\bibitem{barrowtipler} J.~D.~Barrow and F.~J.~Tipler, \emph{The Anthropic Cosmological Principle},
Oxford University Press, 1986.

\bibitem{euler1755} L.~Euler, \emph{Principes g\'en\'eraux du mouvement des fluides},
presented 1755; M\'em.\ Acad.\ Roy.\ Sci.\ Belles-Lettres Berlin \textbf{11} (1757) 274--315.

\bibitem{cfd} J.~Hoffman and C.~Johnson, \emph{Computational Turbulent Incompressible Flow},
Applied Mathematics: Body and Soul, vol.~4, Springer, 2007.

\bibitem{dalembert} J.~Hoffman and C.~Johnson, \emph{Resolution of d'Alembert's paradox},
J.\ Math.\ Fluid Mech.\ \textbf{12} (2010) 321--334.

\bibitem{flight} J.~Hoffman, J.~Jansson and C.~Johnson, \emph{New theory of flight},
J.\ Math.\ Fluid Mech.\ \textbf{18} (2016) 219--241.

\bibitem{drag} J.~Hoffman, \emph{Efficient computation of mean drag for the subcritical flow past a
circular cylinder using general Galerkin G2}, Int.\ J.\ Numer.\ Meth.\ Fluids \textbf{59} (2009)
1241--1258.

\bibitem{parameterfree} J.~Hoffman, J.~Jansson, N.~Jansson and R.~Vilela de Abreu, \emph{Towards a
parameter-free method for high Reynolds number turbulent flow simulation based on adaptive finite
element approximation}, Comput.\ Methods Appl.\ Mech.\ Engrg.\ \textbf{288} (2015) 60--74.

\bibitem{clay} C.~Johnson, \emph{Constructive Existence and Qualitative Uniqueness of Turbulent
Navier--Stokes Solutions}, companion paper.

\bibitem{constructive} C.~Johnson, \emph{Pure existence says nothing, constructive existence says
everything}, Leibniz World of Math, 2026.

\bibitem{realqmbook} C.~Johnson, \emph{Real Quantum Mechanics} (Body\&Soul, Vol.~VI),
\url{https://claes542.github.io/RealMolecule/realquantum3.pdf}.

\bibitem{realqm} C.~Johnson, \emph{Computational Quantum Mechanics as Multiphase 3D Continuum
Mechanics}, in~\cite{realqmbook}.

\bibitem{universality} C.~Johnson, \emph{Universality in Quantum Mechanics: A Claim That Cannot Be
Tested, and One That Can}, companion paper.

\bibitem{atoms} C.~Johnson, \emph{RealQM for Atoms and the Periodic Table}.

\bibitem{chirality} C.~Johnson, \emph{The Tetrahedral Stereocentre from Packing: Chirality's
Geometric Origin without Hybridisation}, companion paper.

\bibitem{gravity} C.~Johnson, \emph{Charge Below, Gravity Above}, in~\cite{realqmbook}.

\bibitem{education} C.~Johnson, \emph{Mathematics Education when Formal Mathematics is Automated},
companion paper.

\bibitem{ainumerical} C.~Johnson, \emph{AI versus Numerical Analysis: Two Forms of Computation},
companion paper.

\end{thebibliography}

\end{document}
